\section{Variables latentes y Autoencoder}

\subsection{¿Qué es una variable latente (Latent Variable)?}

Para entender el concepto de variable latente, el curso recurre a la célebre \textbf{alegoría de la caverna} (Allegory of the Cave) de la República de Platón: un grupo de prisioneros está atado dentro de una caverna, de frente a una pared, y lo único que pueden ver son las sombras que la luz del fuego proyecta sobre esa pared a partir de los objetos. Para los prisioneros, la sombra es su realidad, pero la estructura real del objeto —la parte que nunca pueden observar directamente— es la causa fundamental que produce esas sombras.

\begin{figure}[H]
\centering
\includegraphics[width=0.85\textwidth]{slides-images/slide-10.jpg}
\caption{La alegoría de la caverna de Platón como analogía del concepto de variable latente\protect\footnotemark}
\end{figure}
\footnotetext{Slides página 10. Fuente: la República de Platón.}

\begin{knowledgebox}{Analogía filosófica de la variable latente}
En el aprendizaje automático, los datos que observamos son como las sombras en la pared: son la manifestación de ciertos factores subyacentes (variables latentes). Las variables latentes (Latent Variables) son como los objetos reales que proyectan las sombras en la caverna; aunque no se pueden medir directamente, determinan la distribución de los datos que observamos. Uno de los objetivos centrales del modelado generativo es extraer estas variables latentes a partir de los datos observados.
\end{knowledgebox}

\subsection{Principio básico del Autoencoder}

El \textbf{Autoencoder} (autocodificador) es un método de aprendizaje no supervisado que se usa para aprender representaciones de características de baja dimensión a partir de datos de entrenamiento sin etiquetar. Su idea central es: mediante un codificador (Encoder) se comprime el dato de entrada de alta dimensión a un espacio latente (Latent Space) de baja dimensión, y luego un decodificador (Decoder) intenta reconstruir el dato de entrada original a partir del espacio latente.

\begin{figure}[H]
\centering
\includegraphics[width=0.85\textwidth]{slides-images/slide-13.jpg}
\caption{Autoencoder: el codificador mapea el dato $x$ a un espacio latente de baja dimensión $z$\protect\footnotemark}
\end{figure}
\footnotetext{Slides página 13.}

El mapeo que aprende el codificador:

$$\text{Encoder}: x \to z \quad \text{(alta dimensión} \to \text{baja dimensión)}$$

Donde $x$ es el dato de entrada (por ejemplo, una imagen) y $z$ es el vector de variables latentes. La razón de mantener la dimensión de $z$ muy por debajo de la de $x$ es que queremos que la red \textbf{comprima} los datos y extraiga la representación de características más esencial.

\subsection{Aprender el espacio latente mediante la reconstrucción}

Pregunta clave: ¿cómo entrenar al codificador para que aprenda buenas variables latentes $z$ sin contar con etiquetas?

La respuesta es: conectar un \textbf{decodificador} (Decoder) después del codificador y entrenar la red completa para reconstruir el dato de entrada original.

\begin{figure}[H]
\centering
\includegraphics[width=0.85\textwidth]{slides-images/slide-15.jpg}
\caption{Autoencoder completo: codificación, decodificación y reconstrucción\protect\footnotemark}
\end{figure}
\footnotetext{Slides página 15.}

$$\text{Decoder}: z \to \hat{x} \quad \text{(baja dimensión} \to \text{alta dimensión)}$$

El objetivo de entrenamiento es minimizar la distancia entre el dato de entrada $x$ y la salida reconstruida $\hat{x}$:

$$\mathcal{L}(x, \hat{x}) = \|x - \hat{x}\|^2$$

\begin{itemize}
    \item $x$: dato de entrada original
    \item $\hat{x}$: dato reconstruido después de la codificación y la decodificación
    \item $\|\cdot\|^2$: error cuadrático medio (Mean Squared Error)
\end{itemize}

\begin{importantbox}{La naturaleza autosupervisada del Autoencoder}
¡La función de pérdida del Autoencoder \textbf{no necesita ninguna etiqueta externa}! Usa el dato de entrada mismo como señal de aprendizaje y, al minimizar el error de reconstrucción, obliga a la red a aprender una representación comprimida significativa. Esto hace del Autoencoder un método de aprendizaje sin etiquetas \textbf{autosupervisado}.
\end{importantbox}

\subsection{Dimensión del espacio latente y calidad de la reconstrucción}

La dimensión del espacio latente influye directamente en la calidad de la reconstrucción. El autocodificado es, en esencia, una operación de \textbf{compresión}: cuanto más pequeño es el espacio latente, mayor es el cuello de botella de información, y la red se ve obligada a aprender una representación más compacta.

\begin{figure}[H]
\centering
\includegraphics[width=0.85\textwidth]{slides-images/slide-17.jpg}
\caption{Efecto de la dimensión del espacio latente en la calidad de la reconstrucción: 2D vs. 5D vs. dato real\protect\footnotemark}
\end{figure}
\footnotetext{Slides página 17.}

Como se muestra en la figura, cuando el espacio latente es de 2D, el resultado de la reconstrucción es notoriamente borroso; al aumentar a 5D, la calidad de la reconstrucción mejora de manera considerable. Esto refleja el compromiso entre la capacidad de representación y el grado de compresión.

\subsection{Resumen de la sección}

El Autoencoder logra el aprendizaje de características sin etiquetas mediante una arquitectura de codificador-decodificador. Comprime los datos a un espacio latente de baja dimensión y luego reconstruye el dato de entrada original a partir de ese espacio latente. Pero el Autoencoder tradicional presenta un problema clave: el proceso de codificación es \textbf{determinista} (deterministic); dado el mismo dato de entrada, siempre produce la misma variable latente. Esto significa que el espacio latente puede no ser lo bastante suave ni continuo como para usarse eficazmente en la generación de muestras nuevas. Esto da paso a la sección siguiente sobre el Variational Autoencoder.
